(1) 在阱内可分离变量 \(\psi=X(x)Y(y)Z(z)\)，每个坐标边界波函数为零，故归一的一维本征函数为 \(\sqrt{2/L}\sin(n\pi x/L)\)，\(n=1,2,\ldots\)。三维归一波函数是乘积 \(\psi_{n_xn_yn_z}=(2/L)^{3/2}\sin(n_x\pi x/L)\sin(n_y\pi y/L)\sin(n_z\pi z/L)\)。
(2) 分离能量相加给 \(E_{\mathbf n}=\hbar^2\pi^2(n_x^2+n_y^2+n_z^2)/(2mL^2)\)。最小值对应 \((1,1,1)\)，只有一种空间态，故基态不简并。第一激发能级对应 \((2,1,1),(1,2,1),(1,1,2)\)，能量平方数和为6，空间简并度为3；若粒子有自旋，应另外计数。
(3) 基态三个方向相互独立，\(E_{111}=3\pi^2\hbar^2/(2mL^2)\)。由对称性 \(\langle x\rangle=L/2\)；再计算 \(\langle x^2\rangle=(2/L)\int_0^L x^2\sin^2(\pi x/L)dx=L^2(1/3-1/(2\pi^2))\)，同理适用于 \(y,z\)。
(4) 令 \(n_x^2+n_y^2+n_z^2\le R^2\)，其中 \(R^2=2mL^2E/(\pi^2\hbar^2)\)。大量子数时正整数格点数近似正八分球体积 \(N(E)\simeq(\pi/6)R^3=V(2mE)^{3/2}/(6\pi^2\hbar^3)\)，其中 \(V=L^3\)；它是无自旋简并的领先态数。
